%=== CHAPTER THREE (3) ===
%=== (Actual work done and contribution, including literature survey) ===

\chapter{Methodology}

\section{First Section}

Describe your method here. The examples below show the elements used in this template.

\subsection{Lists}

Numbered list:

\begin{enumerate}
    \item \textbf{First step:} Description of the first step.
    \item \textbf{Second step:} Description of the second step.
\end{enumerate}

Bulleted list:

\begin{itemize}
    \item \textbf{First item:} Description of the first item.
    \item \textbf{Second item:} Description of the second item.
\end{itemize}

\subsection{Equations}

Inline maths is written as \(y = ax + b\). A displayed equation without a number:

\[
K =
\begin{bmatrix}
f_x & 0 & c_x \\
0 & f_y & c_y \\
0 & 0 & 1
\end{bmatrix}
\]

A numbered equation, referred to as Equation \eqref{eq3.1}:

\begin{equation}
E = mc^2
\label{eq3.1}
\end{equation}

\subsection{Tables}

Table \ref{table:example} shows the table format. The caption is placed above the table.

\begin{table}[H]
\caption{Table Caption Goes Here}
\centering
\begin{tabular}{cc}
\toprule
\textbf{Column One} & \textbf{Column Two} \\
\midrule
Row 1 & 0.0000 \\
Row 2 & 0.0000 \\
Row 3 & 0.0000 \\
\bottomrule
\end{tabular}
\label{table:example}
\end{table}

\subsection{Code}

Code snippets are placed in a \texttt{verbatim} environment:

\begin{verbatim}
def example_function(x):
    return x + 1
\end{verbatim}

\subsubsection{A Subsubsection}

Headings are numbered down to this level.


%=== END OF CHAPTER THREE ===
\newpage
